\section{Baseline Authenticated Encryption Bounds}
\label{sec:aead-security}

\Cref{sec:root-tag-hash-security} proves the tag-centric properties of \TW{}.
Here we discharge the AEAD-facing obligations: multi-user indistinguishability,
its single-user all-in-one AE corollary, and \textsf{INT-RUP} integrity. These
bounds use the same hidden-key, leaf-collision, and tag-guessing accounting as
the tag analysis, but their role is narrower: they show that the mode remains a
conventional nonce-respecting AEAD.

We first prove mu-ind in the flat random oracle model, lift the result to the
public permutation model using \Cref{app:flat-sponge-lift}, and derive
all-in-one AE as the $D = 1$ case. We then prove \textsf{INT-RUP} in the
separated-interface model of \cite{ABLMMY14}, where plaintext-computing queries
increase the construction work counts but do not change the non-replay
verification count. All concrete instantiations are summarized in
\Cref{sec:concrete-security}.

\subsection{AEAD Public Permutation Games}
\label{sec:aead-games}

The public-permutation games in this section use the primitive, validity, and
construction-oracle conventions of \Cref{sec:games-pp}. The random oracle
versions are obtained by the mechanical substitution of \Cref{sec:ro-model}.

\paragraph{All-in-one AE.}
The all-in-one authenticated encryption game~\cite{RS06} samples
$K \sample \bits^k$ and compares the real pair
$(\Enc_K[\pi], \Ver_K[\pi])$ with the ideal pair $(\Rand, \Replay)$.
On each accepted encryption query $(U, A, M)$, $\Rand$ samples a uniformly
random byte string of length $\|M\| + \tauroot/8$, parses it as
$C \concat T$ with $T \in \bits^{\tauroot}$, records $(U, A, C, T)$, and
returns $C \concat T$; $\Replay$ accepts exactly recorded tuples. Encryption
queries are nonce-respecting, and
\[
  \Advae_\TW(\mathcal{A}) =
  \bigl|
    \Pr[\mathcal{A}^{\Enc_K[\pi],\, \Ver_K[\pi],\, \pi,\, \pi^{-1}} \Rightarrow 1]
    - \Pr[\mathcal{A}^{\Rand,\, \Replay,\, \pi,\, \pi^{-1}} \Rightarrow 1]
  \bigr|.
\]
This all-in-one AE advantage is the single-user case of the multi-user
real-vs-ideal notion of \Cref{sec:mu-ind}; \Cref{cor:ae} derives it from
\Cref{thm:mu-ind}.

\paragraph{RUP integrity.}
The \textsf{INT-RUP} game samples $K \sample \bits^k$ and
$\pi \sample \Perm(\bits^{1600})$, exposes $\Enc_K[\pi]$,
$\PDec_K[\pi]$, $\Ver_K[\pi]$, and $(\pi, \pi^{-1})$, and maintains the replay
table $W$ of encryption outputs. Accepted encryption queries are
nonce-respecting: an encryption query is rejected if its nonce has appeared in a
previous accepted encryption query. Plaintext-computing and verification
queries may use arbitrary nonces. On an accepted encryption query $(U, A, M)$ the
game returns $\Enc_K[\pi](U, A, M)$, parsed and recorded as $(U, A, C, T)$ in $W$.
On a valid plaintext-computing query $(U, A, C, T)$ it returns
$\PDec_K[\pi](U, A, C, T)$, and on an invalid query it returns $\bot$. On a valid
verification query $(U, A, C, T)$ it returns $\Ver_K[\pi](U, A, C, T)$; if this bit
is $1$ and $(U, A, C, T) \notin W$ when the query is asked, the game records a
forgery. Invalid verification queries return $0$. The game returns $1$ iff a
forgery is recorded, and
\[
  \AdvintRUP_\TW(\mathcal{A}) =
  \Pr[\mathsf{G}^{\mathsf{INT\text{-}RUP}}_\TW(\mathcal{A}) \Rightarrow 1].
\]

\subsection{Multi-User Setting}
\label{sec:mu-setting}

This subsection instantiates the mu-ind game of \cite{BT16}
(\Cref{sec:mu-ind}) for \TW{} in the public permutation model, under
the primitive convention of \Cref{sec:games-pp}. The challenger
samples a challenge bit $b \sample \bits$ and a permutation
$\pi \sample \Perm(\bits^{1600})$, and exposes $(\pi, \pi^{-1})$ as the
primitive interface in both worlds. The
adversary $\mathcal{A}$ creates users on demand via $\mathsf{New}$, and
the game maintains an active-user counter $v$. We write $D$ for an
execution-uniform cap on the number of $\mathsf{New}$ calls. Equivalently,
the game may sample independent keys $K_1,\dots,K_D$ at the start of the
experiment and let the $d$-th $\mathsf{New}$ call activate $K_d$; inactive
indices are never valid users.

The mu-ind construction oracles inherit the validity convention of
\Cref{sec:games-pp}. The encryption oracle on input $(d, U, A, M)$
rejects, without recording, if $(U, A, M)$ is outside the \TW{} encryption
domain, if $d \notin \{1, \dots, v\}$, or if $(d, U) \in V$; otherwise
it sets $C = \Enc_{K_d}[\pi](U, A, M)$ when $b = 1$, and samples a
uniformly random byte string $C$ of length $\|M\| + \tauroot/8$ when
$b = 0$. The game records
$(d, U)$ in $V$ and $(d, U, A, C)$ in $W$, and returns $C$. The
verification oracle on input $(d, U, A, C)$ rejects if $d \notin
\{1, \dots, v\}$ or if $(U, A, C)$ is not a valid \TW{}
verification/decryption input; if $(d, U, A, C) \in W$ it returns
$\mathsf{true}$; if $b = 0$ it returns $\mathsf{false}$; otherwise it
returns $[\Dec_{K_d}[\pi](U, A, C) \neq \bot]$. Nonce uniqueness is
enforced per user by the set $V$.

After querying these oracles, $\mathcal{A}$ outputs $b' \in \bits$.
The mu-ind advantage of $\mathcal{A}$ against \TW{} is
\[
  \Advmuind_\TW(\mathcal{A})
  = |2 \Pr[b' = b] - 1|.
\]
A valid verification query is \emph{non-replay} if its tuple is not in
$W$ when the query is asked; replay hits are answered by the recorded
transcript and are not counted in the verification-query resource below.

Resource caps split per potential user $d \leq D$: $N^{(d)}$,
$\qv^{(d)}$, $\nleaf^{(d)}$ are the per-user counts under
\Cref{sec:resources}, with inactive users contributing zero. The global
totals $N = \sum_{d=1}^D N^{(d)}$, $\qv = \sum_{d=1}^D \qv^{(d)}$,
$\nleaf = \sum_{d=1}^D \nleaf^{(d)}$ appear in the final bound: $N$
through the lift term $\Lift_c(\Ntot)$, and $\nleaf$ and $\qv$ through
the leaf-collision and verification-tag terms. $\Nprim$
remains a single shared count over the primitive interface.

The mu-ind treatment applies only to indistinguishability. The
committing security games already give the adversary control of all
$s$ keys, so the $s$-way CMTDs-$\ell$ statements of
\Cref{sec:root-tag-games} are themselves the multi-key committing statement
and need no separate $\mathsf{mu\text{-}cmt}$ notion.

All random oracle adversaries in this section are flat-RO adversaries
(\Cref{sec:ro-model}) with the displayed resources, and the flat sponge lift
transfers each random oracle bound using the derived-adversary cap
$\qRO(\mathcal{B}) \leq \Nprim$.

\begin{theorem}[Random Oracle Multi-User Indistinguishability]
\label{thm:mu-ind-ro}
For every flat-RO mu-ind adversary $\mathcal{B}$ with
execution-uniform user cap $D$ and resource counts of
\Cref{sec:mu-setting},
\[
  \Adv^{\mathsf{mu\text{-}ind}}_{\RO}(\mathcal{B})
  \;\leq\;
  \frac{D(D - 1)}{2^{k + 1}}
  + D \cdot \KeyGuess_k(\qRO(\mathcal{B}))
  + \frac{\nleaf(\mathcal{B})}{2^{\tauleaf}}
  + \frac{\qv(\mathcal{B})}{2^{\tauroot}}.
\]
Here $\qv$ is the non-replay verification-query count of
\Cref{sec:resources}.
\end{theorem}

\begin{proof}
Use the pre-sampled-key formulation of \Cref{sec:mu-setting}: sample
keys $K_1,\dots,K_D$ at the start of the experiment, and let the $d$-th
$\mathsf{New}$ call activate $K_d$, unused keys being ignored.
Let $\mathsf{coll}$ denote the event that two of the pre-sampled keys
$K_1,\dots,K_D$ collide;
$\Pr[\mathsf{coll}] \leq \binom{D}{2}/2^k = D(D-1)/2^{k+1}$ by union
bound. In the random oracle model, use a standard hybrid over the $D$
users: $H_d$ answers users $1,\dots,d$ with $(\Rand,\Replay)$ and
users $d+1,\dots,D$ with the real construction interface, so $H_0$ is
the real mu-ind game and $H_D$ is the ideal one. Let $H_d^\star$ be
$H_d$ with a fixed output bit on $\mathsf{coll}$. Since
$\mathsf{coll}$ has the same probability in all hybrids,
\[
  \Adv^{\mathsf{mu\text{-}ind}}_{\RO}(\mathcal{B})
  \leq
  \Pr[\mathsf{coll}]
  +
  \bigl|\Pr[\mathcal{B}^{H_0^\star} \Rightarrow 1]
       - \Pr[\mathcal{B}^{H_D^\star} \Rightarrow 1]\bigr|.
\]
By the triangle inequality, it remains to bound each adjacent stopped
hybrid. For each $d$, \Cref{lem:mu-adjacent-hybrid} gives
\[
  \bigl|\Pr[\mathcal{B}^{H_{d-1}^\star} \Rightarrow 1]
       - \Pr[\mathcal{B}^{H_d^\star} \Rightarrow 1]\bigr|
  \leq
  \KeyGuess_k(\qRO(\mathcal{B}))
  + \frac{\nleaf^{(d)}}{2^{\tauleaf}}
  + \frac{\qv^{(d)}}{2^{\tauroot}}.
\]
\Cref{lem:mu-adjacent-hybrid} is the only construction-specific step in
the baseline mu-ind proof. It couples the adjacent user-$d$ hybrids through
the all-reject continuation of \Cref{def:adjacent-hybrid-games}. The
coupling agrees until a direct random-oracle query tests the hidden key
$K_d$ or a non-replay verification comparison succeeds. The resulting
first-occurrence decomposition has exactly the three terms displayed above:
the foreign-key test is bounded by
$\KeyGuess_k(\qRO(\mathcal{B}))$; a same-slot collision between final leaf
tags for user $d$ is bounded by
$\nleaf^{(d)}/2^{\tauleaf}$ via \Cref{lem:leaf-collision}; and a successful
verification-first root comparison is bounded by
$\qv^{(d)}/2^{\tauroot}$ via \Cref{lem:ver-first-root-guess}. Encryption
transcript marginal uniformity and environment neutrality
(\Cref{lem:enc-transcript-marginal-uniformity,lem:env-neutrality}) are the
facts that make this coupling independent of the other users.

Summing across $d = 1,\dots,D$ gives
$D \KeyGuess_k(\qRO(\mathcal{B}))$ for the key-test contribution. The
leaf tag contribution sums exactly:
\[
  \sum_d \frac{\nleaf^{(d)}}{2^{\tauleaf}}
  = \frac{\nleaf(\mathcal{B})}{2^{\tauleaf}},
\]
since $\sum_d \nleaf^{(d)} = \nleaf(\mathcal{B})$. The
verification-tag contribution satisfies
\[
  \sum_d \qv^{(d)} / 2^{\tauroot}
  = \qv(\mathcal{B}) / 2^{\tauroot}.
\]
Combining these estimates with the $\mathsf{coll}$ penalty gives the
stated RO bound.
\end{proof}

\begin{theorem}[Multi-User Indistinguishability]
\label{thm:mu-ind}
For every mu-ind adversary $\mathcal{A}$ against \TW{} with
execution-uniform user cap $D$ and the global resources of
\Cref{sec:mu-setting},
\[
  \Advmuind_\TW(\mathcal{A})
  \;\leq\;
  \Lift_c(\Ntot)
  + \Lift_c(\Nprim)
  + \frac{D(D - 1)}{2^{k + 1}}
  + D \cdot \KeyGuess_k(\Nprim)
  + \frac{\Nleaf}{2^{\tauleaf}}
  + \frac{\Qv}{2^{\tauroot}}.
\]
\end{theorem}

\begin{proof}
By the Lift Application corollary (\Cref{cor:lift-application}) for the
mu-ind game---a real-vs-ideal game, so the ideal-side proximity term
applies---there is a random oracle mu-ind adversary
$\mathcal{B} = \mathcal{B}_{\mathrm{derived}}(\mathcal{A})$ with
$\qRO(\mathcal{B}) \leq \Nprim$, $\nleaf(\mathcal{B}) \leq \Nleaf$,
$\qv(\mathcal{B}) \leq \Qv$, and the per-user construction-side caps of
\Cref{sec:mu-setting} such that
\[
  \Advmuind_\TW(\mathcal{A})
  \;\leq\;
  \Lift_c(\Ntot) + \Adv^{\mathsf{mu\text{-}ind}}_{\RO}(\mathcal{B})
  + \Lift_c(\Nprim).
\]
\Cref{thm:mu-ind-ro} bounds the random oracle term. Substituting
$\qRO(\mathcal{B}) \leq \Nprim$,
$\nleaf(\mathcal{B}) \leq \Nleaf$, and
$\qv(\mathcal{B}) \leq \Qv$ gives the stated bound.
\end{proof}

\begin{corollary}[All-in-One AE Security]
\label{cor:ae}
For every nonce-respecting AE adversary $\mathcal{A}$ against \TW{}
with the resources of \Cref{sec:resources},
\[
  \Advae_\TW(\mathcal{A})
  \;\leq\;
  \Lift_c(\Ntot)
  + \Lift_c(\Nprim)
  + \KeyGuess_k(\Nprim)
  + \frac{\Nleaf}{2^{\tauleaf}}
  + \frac{\Qv}{2^{\tauroot}}.
\]
\end{corollary}

\begin{proof}
Build a mu-ind adversary $\mathcal{A}^{\mu}$ that makes one
$\mathsf{New}$ query, runs $\mathcal{A}$, and forwards each
encryption and verification query of $\mathcal{A}$ to user $1$. The
primitive interface is forwarded unchanged. The real and ideal views of
$\mathcal{A}$ in the all-in-one AE experiment are exactly the real and
ideal views of $\mathcal{A}^{\mu}$ in the mu-ind experiment with
$D = 1$, and the resource counts are unchanged. On a replayed
encryption tuple the two real verification oracles agree: the mu-ind
oracle short-circuits through its replay table, while the AE oracle
re-runs decryption, which accepts by \TW{} decryption correctness
(\Cref{sec:dec}); all other queries are answered by identical
code. Applying
\Cref{thm:mu-ind} with $D = 1$ eliminates the key-collision term and
gives the stated bound.
\end{proof}

The \textsf{INT-RUP} bound is AEAD-facing, but it is closer to the tag-hash
analysis than to the mu-ind baseline: plaintext-computing queries expose
unverified stream plaintext, and freshness must still reduce to hidden key
tests, leaf-tag collisions, and root-tag guessing. The separated interface
changes the accounting by adding plaintext-computing work to $N$ and
$\nleaf$, while leaving the non-replay verification count $\qv$ unchanged.

\begin{theorem}[Random Oracle \textsf{INT-RUP} Integrity]
\label{thm:int-rup-ro}
For every flat-RO \textsf{INT-RUP} adversary $\mathcal{B}$ against \TW{}
with nonce-respecting encryption queries and the resources of
\Cref{sec:resources},
\[
  \AdvintRUP_{\RO}(\mathcal{B})
  \;\leq\;
  \KeyGuess_k(\qRO(\mathcal{B}))
  + \frac{\nleaf(\mathcal{B})}{2^{\tauleaf}}
  + \frac{\qv(\mathcal{B})}{2^{\tauroot}}.
\]
Plaintext-computing queries contribute to $N$ and $\nleaf$, but not to
$\qv$.
\end{theorem}

\begin{proof}
The proof uses the same hidden-tag accounting as the adjacent-hybrid bound, with
the plaintext-computing interface included in the construction transcript.
Work in the flat-RO \textsf{INT-RUP} game of \Cref{sec:ro-model}. Let
$\forge$ be the event that a valid non-replay verification query accepts.
Let $\badkey$ be the event that a direct $\ROflat$ query decodes under
$\KeyedTWOut$ to the hidden key $K$ and a counted output-point class, and let
$\badleaf$ be the event that some construction-generated final leaf
transcript receives the same $\tauleaf$-bit leaf tag projection as a
distinct final leaf transcript that shares its key, nonce, and leaf index
and is reached by an encryption computation (the same-slot event of
\Cref{lem:leaf-collision}). Final leaf
transcripts generated by encryption, plaintext-computing, and verification
computations are all included in $\nleaf$.

The key-test term is unchanged by the plaintext-computing oracle. By
\Cref{lem:ver-key-uniformity}, stated for the \textsf{INT-RUP} extension,
the hidden key remains uniform before $\badkey$ outside the candidate keys
already tested by direct queries, and \Cref{lem:key-tests} gives
\[
  \Pr[\badkey] \leq \KeyGuess_k(\qRO(\mathcal{B})).
\]

For the leaf term, no direct query has pre-read a hidden-key final leaf tag
point before $\badkey$: such a query would decode under $\KeyedTWOut$ to $K$
and the corresponding $\Ltag(j)$ class. Coupling the real
\textsf{INT-RUP} execution to the all-reject continuation used in
\Cref{lem:int-rup-root-guess} preserves the construction computations and
leaf-tag table accesses up to the first accepting non-replay submission.
The same-slot case of \Cref{lem:leaf-collision}, applied to the enlarged set
of construction computations, gives
\[
  \Pr[\badleaf \text{ before } \badkey \cup \forge]
  \leq
  \frac{\nleaf(\mathcal{B})}{2^{\tauleaf}}.
\]

It remains to bound accepting verification before $\badkey \cup \badleaf$.
The \textsf{INT-RUP} root-guessing lemma (\Cref{lem:int-rup-root-guess})
gives
\[
  \Pr[\forge \text{ before } \badkey \cup \badleaf]
  \leq \qv(\mathcal{B}) / 2^{\tauroot}.
\]
Combining the three bounds through the first-occurrence decomposition
of $\Pr[\forge]$ proves the theorem.
\end{proof}

\begin{theorem}[\textsf{INT-RUP} Integrity]
\label{thm:int-rup}
For every \textsf{INT-RUP} adversary $\mathcal{A}$ against \TW{} with
nonce-respecting encryption queries and the resources of \Cref{sec:resources},
\[
  \AdvintRUP_\TW(\mathcal{A})
  \;\leq\;
  \Lift_c(\Ntot)
  + \KeyGuess_k(\Nprim)
  + \frac{\Nleaf}{2^{\tauleaf}}
  + \frac{\Qv}{2^{\tauroot}}.
\]
\end{theorem}

\begin{proof}
By the Lift Application corollary (\Cref{cor:lift-application}) for the
\textsf{INT-RUP} game, there is a random oracle \textsf{INT-RUP} adversary
$\mathcal{B} = \mathcal{B}_{\mathrm{derived}}(\mathcal{A})$ with
$\qRO(\mathcal{B}) \leq \Nprim$, $\nleaf(\mathcal{B}) \leq \Nleaf$, and
$\qv(\mathcal{B}) \leq \Qv$ such that
\[
  \AdvintRUP_\TW(\mathcal{A})
  \;\leq\;
  \Lift_c(\Ntot) + \AdvintRUP_{\RO}(\mathcal{B}).
\]
\Cref{thm:int-rup-ro} bounds the random oracle term, and substituting the
derived-adversary resource caps gives the stated bound.
\end{proof}
